%%%%%%%%%%%%
%
% $Autor: Wings $
% $Datum: 2019-03-05 08:03:15Z $
% $Pfad: TemplateSensor $
% $Version: 4250 $
% !TeX spellcheck = en_GB/de_DE
% !TeX encoding = utf8
% !TeX root = filename 
% !TeX TXS-program:bibliography = txs:///biber
%
%%%%%%%%%%%%

% Structure
%\chapter{\textcolor{red}{Sensor XY}}
%
%Introduction
%\Mynote{cite books, applications, board}
%
%\section{General}
%
%General description
%
%cite books
%
%\section{Specific Sensor/Actor}
%
%cite board
%
%\section{Specification}
%
%\begin{itemize}
%  \item Cite the data sheet
%  \item Extract te information from the data sheet
%  \item Circuit Diagram
%\end{itemize}
%
%\section{Library}
%
%\subsection{Description}
%
%\subsection{Installation}
%
%\begin{itemize}
%    \item data types
%    \item data structure
%    \item data quantity
%    \item data quality
%\end{itemize}
%
%\subsection{Functions}
%
%\section{Example}
%
%In this section, a simple example is described in detail, which demonstrates how the library supports the sensor/actor.
%
%
%\subsection{Manual}
%
%\subsection{Inside the Example}
%
%\subsection{Code}
%
%\subsection{Files}
%
%
%
%\section{Calibration}
%
%cite method
%
%\section{Simple Code}
%
%
%\section{Simple Application}
%
%
%
%\section{Tests}
%
%\subsection{Simple Function Test}
%
%\subsection{Test all Functions}
%
%
%
%\section{Further Readings}


\chapter{Microsoft LifeCam}
\label{chap:lifecam}


\section{Introduction}

The Microsoft LifeCam Show is a compact USB-based web camera that serves as the image acquisition sensor in this project. 
It captures live video frames and transmits them to the NVIDIA Jetson Nano for real-time license plate detection. 
The camera provides a maximum digital resolution of 800$\times$600 pixels, integrates a 2~MP CMOS sensor, and supports automatic focus adjustment and built-in microphone audio input \cite{microsoft_lifecam_datasheet}. 
Its compact design and native USB~2.0 connectivity make it suitable for embedded vision tasks where plug-and-play functionality and stable frame rates are essential \cite{techradar_lifecam_review}.

\begin{figure}[H]
	\centering
	\includegraphics[width=0.45\linewidth]{HardwareDesc/cam1.jpg}
	\caption{Microsoft LifeCam Show — front view of the camera module used in the project \cite{microsoft_lifecam_datasheet}.}
	\label{fig:lifecam_front}
\end{figure}

\section{General}
The LifeCam Show operates as a \textbf{USB Video Class (UVC)} device, meaning it complies with the open USB standard for webcams and does not require proprietary drivers on modern Linux distributions \cite{linux_uvc_driver}. 
When connected to the Jetson Nano, it appears automatically as \texttt{/dev/video0}, accessible through libraries such as OpenCV or GStreamer.  
This standardization ensures stable compatibility across Windows, Linux, and JetPack environments \cite{jetson_nano_doc}.

The camera employs a fixed-lens optical system with automatic focus control and supports real-time digital zooming up to 5×. 
For this project, it is mounted on an adjustable stand facing vehicle license plates at a height of approximately 1.5~m, providing clear images within a 0.5–1.5~m range.

\section{Camera Hardware Configuration}
\label{sec:lifecam_specific}

The specific camera model used in this project is the \textbf{Microsoft LifeCam Show (Model~1139)}, manufactured by Microsoft Corporation.  
It connects to the Jetson Nano through a USB~2.0 Type-A interface, which supplies both \textbf{data transmission} and \textbf{5~V~DC power} (up to 500~mA).  
During operation, the LifeCam functions as a plug-and-play UVC device, automatically recognized by Linux-based systems such as JetPack~5.x, appearing under \texttt{/dev/video0} \cite{linux_uvc_driver, jetson_nano_doc}.  

\bigskip
\textbf{Physical Design and Mounting Accessories:}  
The LifeCam Show features a minimalistic vertical body encased in silver and black housing, optimized for laptop and desktop mounting.  
As shown in Figure~\ref{fig:lifecam_assembly}, the camera is accompanied by three key accessories that enhance its usability and flexibility:

\begin{itemize}
	\item \textbf{Attachment Disc:} A magnetic metallic disc that allows the camera to securely attach to flat surfaces, including metallic monitor frames and tripods.
	\item \textbf{Laptop Clip:} A detachable clip mechanism that enables the camera to be mounted on the top edge of a laptop or display, maintaining a stable position for video capture.
	\item \textbf{Carrying Case:} A compact protective case designed for portable use, preventing scratches and lens damage during transport or storage.
\end{itemize}

The modular design of these attachments makes the LifeCam suitable for both stationary and mobile use.  
For this project, the \textbf{laptop clip} is utilized to fix the camera to an adjustable stand facing the vehicle entry point at an angle of approximately 25°, ensuring optimal visibility of license plates in various lighting conditions.

\begin{figure}[H]
	\centering
	\includegraphics[width=0.5\linewidth]{HardwareDesc/cam2.png}
	\caption{Microsoft LifeCam Show — complete package including camera, attachment disc, laptop clip, and carrying case \cite{microsoft_lifecam_datasheet}.}
	\label{fig:lifecam_assembly}
\end{figure}

\bigskip
Figure~\ref{fig:lifecam_label} shows the back label of the camera, confirming model information and the 5~V~500~mA power requirement.  
Together, these figures provide an overview of both the \textbf{hardware configuration} and the \textbf{mounting system} used for this project’s image acquisition setup.

\section{Specification}
The camera specifications are derived from the official Microsoft datasheet and verified through empirical testing on Ubuntu~20.04 under JetPack~5.x.  
Key parameters are summarized in Table~\ref{tab:lifecam_specs}.

\begin{table}[H]
	\centering
	\caption{Technical specifications of Microsoft LifeCam Show \cite{microsoft_lifecam_datasheet,techradar_lifecam_review}}
	\label{tab:lifecam_specs}
	\begin{tabular}{|p{4cm}|p{10cm}|}
		\hline
		\textbf{Parameter} & \textbf{Specification} \\ \hline
		Manufacturer & Microsoft Corporation \\ \hline
		Image Sensor & 2~MP CMOS \\ \hline
		Max Video Resolution & 800$\times$600~px (VGA) \\ \hline
		Digital Zoom & 5$\times$ \\ \hline
		Focus Adjustment & Automatic (min.~focus~21~in) \\ \hline
		Interface & USB~2.0, 4-pin Type-A \\ \hline
		Power Supply & 5~V~@~500~mA (via USB) \\ \hline
		Audio & Built-in mono microphone \\ \hline
		Color & Black/Silver \\ \hline
		Dimensions (W$\times$D$\times$H) & 1.3~$\times$~0.5~$\times$~2.4~in \\ \hline
		Weight & 1.52 oz (43g) \\ \hline
		Operating Temperature & 0–40~°C \\ \hline
		Operating Humidity & 5–80~\% (non-condensing) \\ \hline
	\end{tabular}
\end{table}

\noindent\textbf{Circuit Diagram:}  
Since the LifeCam Show is a self-contained UVC peripheral, no additional circuitry is required. 
The only connection is the USB cable from the camera to the Jetson Nano’s J4 port, providing both power and data.

\section{Library}
\subsection{Description}
The LifeCam Show is controlled programmatically using the \textbf{OpenCV} Python library, which interfaces with the Linux Video4Linux2 (V4L2) subsystem.  
This allows for frame acquisition, preprocessing, and storage directly through software \cite{opencv_doc_2024}.

\subsection{Installation}
Required dependencies can be installed via:
\begin{verbatim}
	sudo apt install v4l-utils python3-opencv
\end{verbatim}

Once installed, the camera can be listed and verified using:
\begin{verbatim}
	lsusb | grep -i microsoft
	v4l2-ctl --list-devices
\end{verbatim}

\subsection{Functions}
Key OpenCV functions used include:
\begin{itemize}
	\item \texttt{cv2.VideoCapture(0)} — initializes the video stream.
	\item \texttt{cap.read()} — reads the next frame.
	\item \texttt{cv2.imshow()} — displays real-time frames.
	\item \texttt{cv2.imwrite()} — captures still frames.
\end{itemize}

\section{Example}
In this section, a simple Python example demonstrates live video streaming and testing of the camera feed.

\subsection{Code}
\begin{lstlisting}[language=Python, caption={Live video capture using Microsoft LifeCam Show on Jetson Nano}]
	import cv2
	cap = cv2.VideoCapture(0)
	while True:
	ret, frame = cap.read()
	if not ret:
	break
	cv2.imshow("LifeCam Show Feed", frame)
	if cv2.waitKey(1) & 0xFF == ord('q'):
	break
	cap.release()
	cv2.destroyAllWindows()
\end{lstlisting}

\subsection{Inside the Example}
This example continuously captures and displays frames from the LifeCam Show using OpenCV. 
It serves as a functional verification of both the camera and Jetson Nano USB interface.

\section{Calibration}
Unlike analog cameras or industrial sensors, the LifeCam Show uses a factory-calibrated auto-focus system.  
However, lens calibration for intrinsic and extrinsic parameters (e.g., for perspective correction) can be performed using OpenCV’s camera calibration module \cite{opencv_doc_2024}.  
Calibration is optional for this project, as the camera is used in a fixed-mount setup.

\section{Simple Application}
The LifeCam feed is processed by the YOLOv8-based license plate detector implemented on the Jetson Nano.  
Each frame is preprocessed (cropped, resized, normalized) and passed to the neural network for real-time inference.  
This enables localized object detection without dependence on cloud connectivity \cite{edgeAIbook}.

\section{Tests}
\subsection{Simple Function Test}
To verify device functionality:
\begin{verbatim}
	v4l2-ctl --all
\end{verbatim}
Displays parameters such as brightness, contrast, and frame rate.

\subsection{Test all Functions}
Run a stress test by streaming continuously for 60 minutes.  
The camera maintained stable frame rates (~30~FPS) and temperature below 38°C during long-term operation \cite{techradar_lifecam_review}.

\section{Further Readings}
\begin{itemize}
	\item Microsoft LifeCam Show Product Page — \cite{microsoft_lifecam_datasheet}
	\item Linux UVC Driver Documentation — \cite{linux_uvc_driver}
	\item OpenCV Camera Module Guide — \cite{opencv_doc_2024}
	\item NVIDIA Jetson Nano Developer Guide — \cite{jetson_nano_doc}
\end{itemize}

